\documentclass[10pt,twocolumn]{article}
\usepackage[utf8]{inputenc}
\usepackage[T1]{fontenc}
\usepackage{amsmath,amssymb}
\usepackage{graphicx}
\usepackage{booktabs}
\usepackage{hyperref}
\usepackage{cite}
\usepackage{geometry}
\geometry{margin=2.0cm}

\title{\textbf{Silent Physical Qubit Fallback and Gate-Fidelity Collapse Under \texttt{optimization\_level=0} on IBM Heron Processors: A Technical Diagnostic}}

\author{
  \textbf{M\'at\'e R\'obert}\thanks{Corresponding author: \texttt{materobi86@gmail.com}. Open to research collaborations and academic partnerships.}\\
  \textit{Independent Quantum Researcher, Silent Coder Lab}\\
  \textit{Budapest, Hungary}
}
\date{November 2026}

\begin{document}
\maketitle

\begin{abstract}
When benchmarking low-level quantum circuits on cloud superconducting processors, researchers routinely disable compiler heuristics via Qiskit's \texttt{optimization\_level=0} to prevent gate cancellations and evaluate raw physical operations. In this technical note, we report an undocumented calibration metadata trap on 156-qubit IBM Heron architectures (\texttt{ibm\_kingston}, \texttt{ibm\_marrakesh})~\cite{heron2024,krantz_guide}: legacy calls to \texttt{backend.properties()} return $-1$ across physical qubit indices. Consequently, transpiling without an explicit \texttt{initial\_layout} silently falls back to physical qubit 0, which typically exhibits elevated error rates or degraded calibration. On \texttt{ibm\_kingston}, executing an identity identity benchmark ($X \cdot X = I$) under default \texttt{opt\_level=0} yielded an apparent fidelity of only $49.32\%$ ($1024$ shots, Job \texttt{db7v9tor1c5c738lf1gg}), indistinguishable from maximal dephasing. In contrast, dynamically identifying calibrated qubits via an auxiliary high-optimization query (\texttt{opt\_level=3}) and forcing physical qubit 91 via \texttt{initial\_layout=[91]} restored the benchmark fidelity to $99.32\%$ ($1017/1024$ counts) and Loschmidt echo recovery ($U \cdot U^{\dagger}$) to $99.41\%$~\cite{loschmidt_gorin2006}. We provide diagnostic code patterns and recommend standard precautions for researchers conducting gate-level quantum hardware benchmarks.
\end{abstract}

\section{Introduction}

Benchmarking unitary fidelity on superconducting transmon processors requires careful isolation of hardware noise from compiler optimizations. In the Qiskit software development kit (SDK)~\cite{cross2018qiskit}, disabling synthesis passes is traditionally achieved by setting \texttt{optimization\_level=0} in the \texttt{transpile()} function. This ensures that gate sequences (such as echo pulses, randomized benchmarking sequences, or dynamical state inversion $U \cdot U^{\dagger}$) are preserved verbatim without algebraic reduction~\cite{loschmidt_gorin2006,knill2008rb}.

However, compilation passes perform two distinct tasks: (1) algorithmic gate reduction, and (2) physical qubit routing/layout selection~\cite{li2019tackling,sivarajah2020tket}. In this note, we document an empirical pitfall where disabling gate reduction on IBM Heron processors inadvertently deactivates intelligent qubit selection, inducing a severe fidelity collapse.

\section{The Heron Metadata Anomaly}

On IBM's 156-qubit Heron r2 and r3 processors~\cite{heron2024}, the legacy \texttt{backend.properties()} API returns sentinel values ($-1$) for physical qubit readout and gate error metrics. When a user submits an unconstrained compilation job:

\begin{verbatim}
qc_transpiled = transpile(
    qc_raw, 
    backend=backend, 
    optimization_level=0
)
\end{verbatim}

The transpiler's layout pass finds no valid qubit quality ranking in the property cache. Rather than raising a warning or querying modern backend target metrics, the compiler silently falls back to a trivial identity permutation, mapping virtual qubit $q_0$ directly to physical qubit index 0~\cite{cross2018qiskit}.

\section{Hardware Measurements on \texttt{ibm\_kingston}}

To isolate the source of error, we evaluated two distinct circuits on the 156-qubit Heron processor \texttt{ibm\_kingston} (1024 shots per trial):
\begin{enumerate}
  \item \textbf{Identity Calibration Test:} Two consecutive NOT gates ($X \cdot X = I$) executed on initial state $|0\rangle$, where ideal outcome probability is $P(|0\rangle) = 1.0$.
  \item \textbf{Unitary Loschmidt Inversion ($U \cdot U^{\dagger}$):} A 10-gate sequence composed of 6 virtual $\text{RZ}$ and 4 physical $\text{SX}$ gates reversing an arbitrary state rotation back to $|0\rangle$.
\end{enumerate}

\begin{table}[htbp]
\centering
\small
\begin{tabular}{lcccc}
\toprule
\textbf{Circuit Benchmark} & \textbf{Layout Mode} & \textbf{Target QPU} & \textbf{$P(|0\rangle)$} & \textbf{Job ID} \\
\midrule
$X \cdot X = I$ (Aer Sim) & Default & Simulator & $100.00\%$ & \texttt{offline} \\
$X \cdot X = I$ (Hardware) & Auto (Qubit 0) & \texttt{ibm\_kingston} & \textbf{49.32\%} & \texttt{db7v9tor1c5c...} \\
$X \cdot X = I$ (Hardware) & Forced Qubit 91 & \texttt{ibm\_kingston} & \textbf{99.32\%} & \texttt{db7v9wfi99s7...} \\
\midrule
$U$ Control (Hardware) & Forced Qubit 91 & \texttt{ibm\_kingston} & $15.33\%$ & \texttt{db7tvaor1c5c...} \\
$U \cdot U^{\dagger}$ Inversion & Auto (Qubit 0) & \texttt{ibm\_kingston} & \textbf{45.31\%} & \texttt{db7talor1c5c...} \\
$U \cdot U^{\dagger}$ Inversion & Forced Qubit 91 & \texttt{ibm\_kingston} & \textbf{99.41\%} & \texttt{db7tc8or1c5c...} \\
\bottomrule
\end{tabular}
\caption{Measured state survival on \texttt{ibm\_kingston} demonstrating the dramatic fidelity drop caused by unguided qubit 0 allocation under \texttt{optimization\_level=0}.}
\label{tab:heron_benchmarks}
\end{table}

As summarized in Table~\ref{tab:heron_benchmarks}, executing on automatically allocated physical qubit 0 resulted in $P(|0\rangle) = 49.32\%$, perfectly mimicking a fully decohered state (a random coin toss). In contrast, executing the exact same gate sequence on calibrated physical qubit 91 yielded $P(|0\rangle) = 99.32\%$, representing a greater than $2\times$ reduction in measured error.

Furthermore, in the 10-gate unitary inversion sequence ($U \cdot U^{\dagger}$), unguided compilation resulted in $45.31\%$, whereas explicit layout pin to qubit 91 restored state fidelity to \textbf{99.41\%} ($1018$ of $1024$ shots in state $|0\rangle$), achieving a \textbf{6.48$\times$ fidelity improvement} over the control baseline ($15.33\%$).

\section{Diagnostic and Remediation Protocol}

To avoid this pitfall while strictly preserving gate sequences, we recommend a two-pass compilation workflow. In Phase 1, an auxiliary compilation at \texttt{optimization\_level=3} queries the backend's active coupling map and target constraints to extract an optimal physical layout. In Phase 2, the extracted layout is explicitly injected into the unoptimized transpilation~\cite{sivarajah2020tket,wille2019mapping}:

\begin{verbatim}
# Phase 1: Query optimal layout from compiler
aux = transpile(qc, backend=backend, 
                optimization_level=3)
optimal_layout = aux.layout.initial_layout

# Phase 2: Transpile with zero gate cancellation
qc_bench = transpile(
    qc, 
    backend=backend, 
    optimization_level=0, 
    initial_layout=optimal_layout
)
\end{verbatim}

\section{Conclusion}

Disabling compiler optimizations is necessary for hardware characterization, but on commercial architectures with fragmented or evolving calibration metadata, \texttt{optimization\_level=0} can inadvertently deactivate essential layout routing. Researchers observing anomalously low fidelities ($\sim 50\%$) on superconducting processors should inspect physical qubit assignments before attributing losses to intrinsic dephasing~\cite{krantz_guide,blais_cqed}.

\section*{Data Availability}
Job logs and verification scripts are available at: \url{https://github.com/silentnoisehun/quantum-anchored-memory}.


\begin{thebibliography}{99}
\bibitem{heron2024}
IBM Quantum,
``IBM Quantum Heron processor architecture and calibration specifications,''
IBM Technical Whitepaper, 2024.

\bibitem{cross2018qiskit}
A.~Cross,
``The IBM Q experience and Qiskit open-source quantum computing software,''
\emph{APS March Meeting Abstracts}, vol.~2018, pp.~L58--003, 2018.

\bibitem{krantz_guide}
P.~Krantz, M.~Kjaergaard, F.~Yan, T.~P.~Orlando, S.~Gustavsson, and W.~D.~Oliver,
``A quantum engineer's guide to superconducting qubits,''
\emph{Applied Physics Reviews}, vol.~6, no.~2, p.~021318, 2019.

\bibitem{blais_cqed}
A.~Blais, A.~L.~Grimsmo, S.~M.~Girvin, and A.~Wallraff,
``Circuit quantum electrodynamics,''
\emph{Reviews of Modern Physics}, vol.~93, no.~2, p.~025005, 2021.

\bibitem{li2019tackling}
G.~Li, Y.~Ding, and Y.~Xie,
``Tackling the qubit mapping problem for NISQ-era quantum devices,''
\emph{Proceedings of the International Conference on Architectural Support for Programming Languages and Operating Systems (ASPLOS)}, pp.~1001--1014, 2019.

\bibitem{sivarajah2020tket}
S.~Sivarajah, S.~Dilkes, A.~Cowtan, W.~Howe, and R.~Duncan,
``t|ket>: a retargetable compiler for NISQ devices,''
\emph{Quantum Science and Technology}, vol.~6, no.~1, p.~014003, 2020.

\bibitem{wille2019mapping}
R.~Wille, L.~Burgholzer, and A.~Zulehner,
``Mapping quantum circuits to IBM QX architectures using the minimal number of SWAP and H operations,''
\emph{Proceedings of the 56th Annual Design Automation Conference (DAC)}, pp.~1--6, 2019.

\bibitem{loschmidt_gorin2006}
T.~Gorin, T.~Prosen, T.~H.~Seligman, and M.~\v{Z}nidari\v{c},
``Dynamics of Loschmidt echoes and fidelity decay,''
\emph{Physics Reports}, vol.~435, no.~2--5, pp.~33--156, 2006.

\bibitem{knill2008rb}
E.~Knill et~al.,
``Randomized benchmarking of quantum gates,''
\emph{Physical Review A}, vol.~77, no.~1, p.~012307, 2008.
\end{thebibliography}

\end{document}
